\documentclass[11pt,a4paper]{article}
\usepackage{geometry}
\usepackage{amsmath, amssymb, amsthm}
\usepackage{booktabs}
\usepackage{caption}
\usepackage{subcaption}
\usepackage{float}
\usepackage{hyperref}
\usepackage[authoryear]{natbib}
\usepackage{multirow}
\usepackage{sectsty}
\usepackage{titling}
\usepackage{graphicx}
\usepackage{longtable}

% Academic Journal Formatting
\allsectionsfont{\normalfont\large}
\subsectionfont{\normalfont\normalsize}
\pretitle{\begin{center}\huge\textnormal}
\posttitle{\par\end{center}\vskip 0.5em}
\preauthor{\begin{center}\large\textnormal}
\postauthor{\par\end{center}}
\predate{\begin{center}\large\textnormal}
\postdate{\par\end{center}}

\geometry{margin=1in}

\newcommand{\rsquare}{\ensuremath{\mathcal{R}^2_{\mathrm{C}}}}

% Graphics path for figures
\graphicspath{{causal-targeting-billpay/figs/fausebal/}}

\title{Stable Discovery of Heterogeneous Treatment Effects in the Turkish Bank Overdraft Experiment: \\ A StaDISC Analysis of Bill Payment Messaging}
\author{Technical Analysis Report}
\date{December 28, 2025}

\begin{document}
	
\maketitle
	
\section*{Abstract}
This report presents a comprehensive application of the Stable Discovery of Interpretable Subgroups via Calibration (StaDISC) methodology to the field experiment data from \citet{alan2018} concerning bank overdraft behavior in Turkey. Using the Predictability, Computability, and Stability (PCS) framework developed by \citet{dwivedi2020}, we analyze heterogeneous treatment effects of bill payment messaging on customer financial behavior. From an initial sample of 108,000 customers, we isolate 36,024 customers receiving bill payment promotions and train 23 CATE estimators using 4-fold cross-validation with three random perturbations. Our analysis reveals that while global CATE estimation exhibits poor calibration with negative $\rsquare$-scores on validation folds, local quantile-based subgroups demonstrate stable and statistically significant treatment effects. The top-performing estimator (R-learner with Random Forest outcome model and Lasso effect learner) identifies a subgroup comprising the top 10\% of predicted treatment effects that exhibits a pooled treatment effect of 0.045 with $t$-statistic of 2.90 ($p < 0.01$). This analysis follows the StaDISC methodology through Stage 2 (stability-driven ranking), stopping before the CellSearch procedure for interpretable subgroup characterization.

\section{Introduction}

\subsection{Background and Motivation}
The field experiment conducted by \citet{alan2018} at Yapi Kredi, a major Turkish bank, represents one of the most comprehensive tests of the ``shrouded attributes'' hypothesis in consumer finance. The study sent SMS promotions to 108,000 existing checking account holders, offering discounts on overdraft interest rates along with cross-promotions for automatic bill payment and debit card usage. The counter-intuitive finding---that explicit discount messages \emph{reduced} overdraft usage---supported the theory that unshrouding hidden fees makes them salient, prompting behavioral change.

However, the aggregate Average Treatment Effect (ATE) obscures potentially important heterogeneity. Different customer segments may respond differently to the same intervention based on their financial sophistication, account usage patterns, and pre-existing commitments to automated services. Understanding this heterogeneity is crucial for both theoretical advancement and practical policy design.

\subsection{The StaDISC Framework}
We apply the StaDISC methodology introduced by \citet{dwivedi2020} during their re-analysis of the VIGOR clinical trial data. StaDISC builds on the PCS framework for veridical data science \citep{yu2019}, which emphasizes three core principles. First, \emph{Predictability} demands that models demonstrate out-of-sample prediction accuracy before their conclusions can be trusted. Second, \emph{Computability} requires that analyses be computationally feasible and reproducible. Third, and most distinctively, \emph{Stability} insists that conclusions be robust to reasonable perturbations in data processing, model specification, and random sampling.

The StaDISC workflow proceeds in three stages. The first stage conducts a predictive reality check via calibration, assessing whether CATE estimators can accurately predict treatment effect heterogeneity on held-out data. The second stage performs stability-driven ranking and aggregation of CATE estimators, selecting those that perform consistently across multiple data perturbations. The third stage, CellSearch, translates the quantile-based subgroups into interpretable feature combinations. Our analysis implements the first two stages, identifying stable quantile-based subgroups with significant treatment effects, but does not proceed to the final interpretable cell search step.

\section{Data and Experimental Design}

\subsection{Sample Selection}
The original experiment enrolled 108,000 checking account holders at Yapi Kredi Bank. For this analysis, we focus on a specific treatment arm consisting of customers who received bill payment promotion messages. After applying inclusion criteria requiring the message indicator to equal one and filtering to create clean treatment-control comparisons, our analysis sample contains 36,024 customers. We reserved 20\% of this sample (7,205 customers) as a holdout test set that remained untouched until final validation, leaving 28,819 customers for training and cross-validation. The treatment assignment is approximately balanced, with a propensity score of 0.499.

The treatment variable indicates whether a customer received the bill payment feature activation message combined with the overdraft promotion. Customers in the control group received only the base messaging without the bill payment cross-sell, allowing us to isolate the effect of the bill payment promotion component.

\subsection{Outcome Variable}
The primary outcome measures changes in customer financial behavior following the intervention. Across the full sample, the observed Average Treatment Effect is $\widehat{\text{ATE}} = -0.0158$, indicating a small average reduction in the outcome metric for treated customers relative to controls. This modest effect size presents a challenging setting for heterogeneous treatment effect estimation, as the signal-to-noise ratio is low and detecting systematic heterogeneity requires either larger samples or stronger effect modifiers.

\subsection{Feature Engineering and Selection}
The initial feature space contained 519 variables spanning experimental design strata, categorical indicators for reminder frequency and campaign type, prior feature activations, and numeric variables capturing assets, deposits, payment means, debt levels, minimum balance, and credit card usage patterns. Given this high dimensionality relative to the weak outcome signal, we employed Lasso-based feature selection to reduce overfitting risk.

\begin{figure}[H]
	\centering
	\includegraphics[width=0.95\textwidth]{fausebal_lasso_coef.pdf}
	\caption{Lasso coefficient magnitudes for feature selection. The left panel shows the top 10 features ranked by absolute coefficient value for the treated group, while the right panel shows the corresponding ranking for the control group. Features appearing in either top-10 list were retained for CATE estimation.}
	\label{fig:lasso_selection}
\end{figure}

We fit separate Lasso models for the treated and control groups, selecting the union of the top 10 features from each model based on absolute coefficient magnitude (Figure~\ref{fig:lasso_selection}). This procedure yielded 14 selected features for CATE estimation, striking a balance between model expressiveness and the risk of overfitting given the relatively weak signal strength in the data.

\section{CATE Estimator Library}

\subsection{Estimator Specifications}
Following the StaDISC methodology, we trained a library of 23 CATE estimators spanning four meta-learner frameworks and tree-based methods. The S-learner approach, which fits a single model with treatment as a feature, was implemented with Random Forest and XGBoost base learners. The T-learner strategy, which fits separate outcome models for treated and control groups, employed Lasso, Logistic Regression, Random Forest, and XGBoost. The X-learner framework, designed to handle imbalanced treatment groups, used these same four base learners for the outcome stage with Lasso for the effect stage. Finally, the R-learner, which targets the treatment effect directly through a residual-on-residual regression, was implemented with nine different combinations of Lasso, Random Forest, and XGBoost for the outcome and effect components.

We also trained Causal Tree and Causal Forest estimators, though these exhibited poor stability across perturbations and are consequently excluded from the top rankings discussed below.

All estimators were tuned via 4-fold stratified cross-validation using randomized grid search with 200 iterations per estimator. The hyperparameter grids encompassed regularization strength for Lasso models, tree depth and minimum samples per leaf for Random Forest, and learning rate, subsample ratio, and regularization parameters for XGBoost.

\subsection{Cross-Validation Structure}
A key innovation of StaDISC is its emphasis on stability across data perturbations. For each estimator, we computed CATE predictions under three distinct cross-validation splits: the original 4-fold split used for hyperparameter tuning (random seed 405), and two alternative splits with different random seeds (0 and 7). This design yields 12 training-validation fold pairs per estimator, enabling robust assessment of both predictive performance and stability. Importantly, the tuned hyperparameters from the original split were applied to all three CV configurations, ensuring that any observed instability reflects genuine sensitivity to the data partition rather than variation in model specification.

\section{Calibration-Based Predictive Check}

\subsection{The $\rsquare$ Score}
Following \citet{dwivedi2020}, we assess CATE estimator performance using a calibration-based pseudo-$R^2$ score. For an estimator $\mathfrak{M}$, we partition the feature space into $K=5$ quantile-based bins using the predicted CATE distribution on training folds. Within each bin $\mathcal{G}_j$, we compute both the mean predicted CATE ($\overline{\mathfrak{M}}_{\mathcal{G}_j}$) and the Neyman difference-in-means estimate of the subgroup CATE ($\widehat{\tau}_{\mathcal{G}_j}$). The calibration score measures the weighted average discrepancy between these quantities:
\begin{align}
	\text{Cal-Score}(\mathcal{S}; \mathfrak{M}) &= \sum_{j=1}^{K} \frac{|\mathcal{G}_j \cap \mathcal{S}|}{|\mathcal{S}|} \cdot \left| \overline{\mathfrak{M}}_{\mathcal{G}_j \cap \mathcal{S}} - \widehat{\tau}_{\mathcal{G}_j \cap \mathcal{S}} \right|
\end{align}
We then normalize this score relative to a baseline that uses the constant ATE predictor, yielding the pseudo-$R^2$:
\begin{align}
	\rsquare(\mathcal{S}; \mathfrak{M}) &= 1 - \frac{\text{Cal-Score}(\mathcal{S}; \mathfrak{M})}{\text{Cal-Score}(\mathcal{S}; \widehat{\text{ATE}})}
\end{align}
A score close to 1 indicates excellent calibration, meaning the model's predicted heterogeneity closely matches the observed heterogeneity. Scores near 0 indicate the model performs no better than a constant ATE predictor, while negative scores indicate the model performs worse than this naive baseline.

\subsection{Results: Poor Global Calibration}
Figure~\ref{fig:calibration} presents calibration plots for two representative estimators---the X-learner with Random Forest and the Causal Forest---showing predicted versus observed treatment effects across the five quantile-based bins. The triangular markers represent the mean model predictions within each bin, while the circular markers show the Neyman difference-in-means estimates. The horizontal dashed line indicates the population ATE of $-0.0158$.

\begin{figure}[H]
	\centering
	\includegraphics[width=0.95\textwidth]{fausebal_calibration_plot.pdf}
	\caption{Calibration plots for X-learner (RF) and Causal Forest estimators. Each panel shows predicted CATE (triangles) versus Neyman CATE estimates (circles) across five quantile-based bins, with error bars indicating standard errors. The left column displays results on training folds while the right column shows validation fold performance. The horizontal dashed line marks the population ATE.}
	\label{fig:calibration}
\end{figure}

On the training folds (left panels), both estimators exhibit reasonable calibration---the predicted and observed values track each other across bins. However, on the validation folds (right panels), this correspondence breaks down dramatically. The observed treatment effects show high variance and often fail to follow the monotonic ordering implied by the model predictions.

\begin{figure}[H]
	\centering
	\includegraphics[width=0.85\textwidth]{fausebal_r2_scatter.pdf}
	\caption{Scatter plot of $\rsquare$-scores on training folds (x-axis) versus validation folds (y-axis) for five representative estimators across the four folds of the original CV split. Points below the diagonal indicate overfitting. The left panel shows full $\rsquare$ scores while the right panel shows scores restricted to the bottom two quantile bins.}
	\label{fig:r2_scatter}
\end{figure}

Figure~\ref{fig:r2_scatter} provides a more systematic view, plotting training versus validation $\rsquare$-scores for five representative estimators across the four folds of the original CV split. Points falling below the diagonal indicate overfitting---the model explains heterogeneity on training data that does not generalize to held-out observations. The scatter reveals a clear pattern: training scores cluster in the positive range (0.1--0.5), while validation scores spread widely into negative territory.

\begin{figure}[H]
	\centering
	\includegraphics[width=0.95\textwidth]{fausebal_r2_distribution.pdf}
	\caption{Histogram of $\rsquare$-scores across all 23 estimators and 12 fold combinations (4 folds $\times$ 3 CV perturbations). The left panel shows training fold scores, which concentrate around positive values with mean 0.15. The right panel shows validation fold scores, which center near zero with mean $-0.12$, indicating systematic overfitting.}
	\label{fig:r2_hist}
\end{figure}

Aggregating across all 23 estimators and 12 fold combinations, Figure~\ref{fig:r2_hist} displays the distribution of $\rsquare$-scores. The training fold distribution (left panel) concentrates around positive values with a mean of 0.15 and median of 0.18. In stark contrast, the validation fold distribution (right panel) centers near zero with a mean of $-0.12$ and median of $-0.05$. This systematic gap between training and validation performance is the hallmark of overfitting: the CATE estimators are capturing noise in the training data rather than true treatment effect heterogeneity.

This finding aligns with the original StaDISC analysis of the VIGOR clinical trial data, where similar poor global calibration was observed despite using sophisticated machine learning methods. The poor performance is not unexpected given the signal-to-noise ratio in our setting: with a small ATE magnitude ($-0.0158$) and substantial outcome variance, detecting systematic heterogeneity requires either larger samples or stronger effect modifiers than those available in this dataset.

\section{Local Calibration and Monotonicity Analysis}

\subsection{Quantile-Based Subgroup Generalization}
Despite the disappointing global calibration results, the StaDISC framework does not conclude that heterogeneous treatment effects are entirely absent. Instead, it searches for local pockets of stability by examining whether the ordering of predicted treatment effects generalizes to validation data. The key insight is that even when a CATE estimator fails to predict the precise magnitude of treatment effects, it may still correctly identify which individuals experience larger effects than others.

For a given fraction $q$, we define the quantile-based top subgroup as $\widetilde{\mathcal{G}}_q = \{x : \mathfrak{M}(x) \geq \mathfrak{m}_{1-q}\}$, where $\mathfrak{m}_{1-q}$ is the $(1-q)$-th quantile of predicted CATE values computed on training folds. This subgroup contains the top $q$ fraction of individuals ranked by predicted treatment effect. We then test whether the subgroup treatment effect exceeds that of its complement on the held-out validation fold, computing the indicator $B_q = \mathbf{1}(\widehat{\tau}_{\widetilde{\mathcal{G}}_q \cap \mathcal{V}} \geq \widehat{\tau}_{\widetilde{\mathcal{G}}_q^c \cap \mathcal{V}})$ for subgroup sizes $q \in \{0.1, 0.2, 0.3, 0.4, 0.5\}$.

\subsection{Monotonicity Results}
We also examine pairwise monotonicity---whether adjacent quantile bins maintain their expected ordering on validation data. For instance, we check whether individuals in the [0.8, 0.9] quantile bin exhibit smaller treatment effects than those in the [0.9, 1.0] bin when evaluated on held-out observations.

\begin{figure}[H]
	\centering
	\includegraphics[width=0.95\textwidth]{fausebal_monotonicity_box_plot.pdf}
	\caption{Box plots of pairwise monotonicity generalization frequencies across estimators. Each box shows the distribution of the fraction of validation folds on which the expected ordering holds. The rightmost comparison ([0.8,0.9] vs [0.9,1.0]) shows the strongest generalization, with median frequency around 0.75 for the top estimators.}
	\label{fig:monotonicity}
\end{figure}

Figure~\ref{fig:monotonicity} displays the distribution of monotonicity generalization frequencies across estimators. The results reveal a striking pattern: while middle-range comparisons generalize poorly (median frequencies near 0.5, consistent with random chance), the comparison between the highest bins---[0.8, 0.9] versus [0.9, 1.0]---achieves substantially better generalization. For the top-performing estimator (R-learner with Random Forest outcome and Lasso effect), this comparison generalizes correctly on 75\% of validation folds.

\begin{figure}[H]
	\centering
	\includegraphics[width=0.85\textwidth]{fausebal_new_monotonicity_box_plot.pdf}
	\caption{Box plots comparing quantile-based top subgroups against their complements. Each box shows the distribution across estimators of the frequency with which the top subgroup (defined by threshold $q$) exhibits larger treatment effects than the remaining population on validation folds. The $q=0.1$ threshold (top 10\%) shows the most reliable generalization.}
	\label{fig:top_subgroup_comparison}
\end{figure}

Figure~\ref{fig:top_subgroup_comparison} examines a related but distinct question: how often does the quantile-based top subgroup outperform its complement? For increasingly restrictive thresholds (smaller $q$ values, corresponding to more extreme subgroups), the generalization frequency improves. At the $q = 0.1$ threshold, corresponding to the top 10\% of predicted effects, the top subgroup reliably exhibits larger treatment effects than the remaining 90\% of the population. This finding suggests that while CATE estimators cannot accurately rank individuals across the full distribution, they can successfully identify an extreme subgroup with genuinely elevated treatment effects.

\section{Stability-Driven Ranking of CATE Estimators}

\subsection{$t$-Statistic Computation}
Having established that quantile-based top subgroups generalize better than global CATE predictions, we now turn to the question of estimator selection. The StaDISC framework ranks estimators by the statistical significance of their identified subgroups, measured via $t$-statistics. For a subgroup $\mathcal{G}$, we compute:
\[
T_{\mathcal{G}} = \frac{\widehat{\tau}_{\mathcal{G}} - \widehat{\text{ATE}}}{\sqrt{\widehat{\text{Var}}(\widehat{\tau}_{\mathcal{G}} - \widehat{\text{ATE}})}}
\]
This statistic measures whether the subgroup treatment effect significantly exceeds the population average, with larger positive values indicating stronger evidence for heterogeneity. We aggregate $t$-statistics across quantile thresholds $q \in \{0.1, 0.2, 0.3\}$ and all three CV perturbations, yielding a comprehensive measure of each estimator's ability to identify high-effect subgroups.

\subsection{Estimator Rankings}
Figure~\ref{fig:t_ranking} presents the distribution of mean $t$-statistics across perturbations for each estimator, with estimators sorted by their worst-case rank. The left panel shows the rank distribution while the right panel displays the actual $t$-statistic values. An ideal estimator would show both high $t$-statistics (right panel) and consistent rankings (narrow box in left panel).

\begin{figure}[H]
	\centering
	\includegraphics[width=0.95\textwidth]{fausebal_scores_rank_estimators.pdf}
	\caption{Estimator rankings by $t$-statistic performance across CV perturbations. The left panel shows the distribution of ranks (1 = best) across the three perturbation conditions. The right panel displays the corresponding $t$-statistic values. Estimators are sorted from top to bottom by their worst-case rank, with the most stable performers at the top.}
	\label{fig:t_ranking}
\end{figure}

The rankings reveal substantial heterogeneity in estimator performance. Several estimators achieve high $t$-statistics on some perturbations but collapse on others---for instance, the X-learner with Random Forest achieves the highest single $t$-statistic (0.91 on the original split) but drops to only 0.19 on the second perturbation. This instability disqualifies such estimators under the StaDISC criterion.

Table~\ref{tab:rankings} presents the numerical values underlying Figure~\ref{fig:t_ranking}. The R-learner with Random Forest outcome model and Lasso effect learner (\texttt{r\_rflasso}) emerges as the clear winner, achieving consistently high $t$-statistics across all three perturbations: 0.82 on the original split, 0.84 on the first perturbation, and 0.50 on the second. While this last value is lower than its competitors' peaks, no other estimator maintains such consistent performance across all conditions.

\begin{table}[H]
	\centering
	\caption{Mean $t$-Statistics by Estimator and Perturbation (Top 10 Estimators)}
	\label{tab:rankings}
	\begin{tabular}{lcccc}
		\toprule
		\textbf{Estimator} & \textbf{pert\_none} & \textbf{pert\_cv\_0} & \textbf{pert\_cv\_1} & \textbf{Mean} \\
		\midrule
		\texttt{r\_rflasso} & 0.82 & 0.84 & 0.50 & \textbf{0.72} \\
		\texttt{r\_xgblasso} & 0.80 & 0.69 & 0.35 & 0.62 \\
		\texttt{x\_rf} & 0.91 & 0.75 & 0.19 & 0.61 \\
		\texttt{x\_logistic} & 0.80 & 0.51 & 0.50 & 0.60 \\
		\texttt{t\_lasso} & 0.75 & 0.58 & 0.45 & 0.59 \\
		\texttt{x\_lasso} & 0.71 & 0.56 & 0.46 & 0.57 \\
		\texttt{r\_lassolasso} & 0.72 & 0.50 & 0.44 & 0.55 \\
		\texttt{t\_logistic} & 0.58 & 0.37 & 0.47 & 0.47 \\
		\texttt{t\_xgb} & 0.35 & 0.69 & 0.34 & 0.46 \\
		\texttt{x\_xgb} & 0.58 & 0.26 & 0.08 & 0.30 \\
		\bottomrule
	\end{tabular}
\end{table}

\subsection{Top Estimator Selection}
Following the StaDISC protocol, we select estimators whose rank falls within the top tier across all perturbations. Applying a threshold requiring rank $\leq 7$ on every perturbation, only a single estimator qualifies: the R-learner with Random Forest outcome model and Lasso effect learner (\texttt{r\_rflasso}). This hybrid architecture appears to combine the flexibility of Random Forest for modeling the outcome surface with the regularization of Lasso for the treatment effect, yielding stable performance even when the data partition changes.

\subsection{Stability of Subgroup Membership}
Beyond ranking stability, we assess whether the same individuals are consistently assigned to the top subgroup across different model fits. This membership stability is crucial for practical applications: if different training folds identify entirely different individuals as ``high responders,'' the subgroup definition lacks reliability even if the aggregate statistics generalize.

\begin{figure}[H]
	\centering
	\includegraphics[width=0.85\textwidth]{fausebal_overlap_distribution.pdf}
	\caption{Distribution of pairwise overlap scores for the top 20\% quantile-based subgroup across different CV folds. Each histogram shows, for a given estimator, what fraction of individuals identified as top responders on one fold are also identified on other folds. Higher overlap indicates more stable subgroup definitions.}
	\label{fig:overlap}
\end{figure}

Figure~\ref{fig:overlap} displays the distribution of pairwise overlap scores across estimators. For the \texttt{r\_rflasso} estimator, approximately 80\% of individuals assigned to the top 20\% subgroup on one fold also appear in the top 20\% on other folds. When we relax the threshold to the top 30\%, this overlap increases to 87\%. These high overlap values indicate that the identified subgroup represents a genuine, stable cluster of high-responders rather than an artifact of particular data partitions.

\section{Final Subgroup Validation}

The preceding analyses have identified the R-learner with Random Forest outcome model and Lasso effect learner as the most stable CATE estimator, and have demonstrated that its quantile-based top subgroups generalize reliably to validation data. We now conduct the final validation step: fitting the selected estimator on the full training/validation sample, identifying the top 10\% subgroup, and evaluating its treatment effect on the previously untouched holdout test set.

\subsection{Training/Validation Results}
We first fit the \texttt{r\_rflasso} estimator on all 28,819 customers in the training/validation sample and compute predicted treatment effects. The top 10\% quantile-based subgroup, defined as individuals with predicted effects above the 90th percentile, contains 2,882 customers.

Within this subgroup, the Neyman difference-in-means estimator yields an average treatment effect of 0.0426 with a standard error of 0.0173. The corresponding $t$-statistic of 2.46 indicates statistical significance at the 5\% level (one-sided $p \approx 0.007$). Notably, this subgroup effect is approximately 2.7 times larger in absolute magnitude than the population ATE of $-0.0158$, and it has the opposite sign---suggesting that while the average customer shows a small negative response to the intervention, the top 10\% shows a substantial positive response.

\subsection{Holdout Test Set Validation}
The critical test of our findings comes from the holdout sample, which was reserved at the outset and never used for model training, hyperparameter tuning, or estimator selection. We apply the fitted \texttt{r\_rflasso} model to predict treatment effects for the 7,205 holdout customers, then identify those whose predicted effects exceed the training-set 90th percentile threshold.

This procedure assigns 731 holdout customers (10.1\% of the holdout sample) to the top subgroup. Strikingly, the holdout subgroup exhibits an even larger treatment effect of 0.0538, exceeding the training estimate of 0.0426. The standard error of 0.0340 is larger due to the smaller sample size, yielding a $t$-statistic of 1.58 (one-sided $p \approx 0.057$). While this falls just short of conventional significance thresholds, the consistency in effect magnitude---and the fact that the holdout estimate exceeds the training estimate---provides strong evidence that the identified subgroup captures genuine heterogeneity rather than statistical artifacts.

\subsection{Pooled Analysis}
Combining the training/validation and holdout samples yields the maximum statistical power for estimating the subgroup treatment effect. The pooled sample contains 36,024 customers, with 3,613 (10.0\%) assigned to the top subgroup based on the fitted model's predictions.

\begin{table}[H]
	\centering
	\caption{Summary of Final Subgroup Validation Results}
	\label{tab:final_results}
	\begin{tabular}{lccc}
		\toprule
		& \textbf{Training/Val} & \textbf{Holdout} & \textbf{Pooled} \\
		\midrule
		Total sample size & 28,819 & 7,205 & 36,024 \\
		Subgroup size & 2,882 (10.0\%) & 731 (10.1\%) & 3,613 (10.0\%) \\
		Subgroup ATE & 0.0426 & 0.0538 & \textbf{0.0448} \\
		Standard Error & 0.0173 & 0.0340 & 0.0154 \\
		$t$-statistic & 2.46 & 1.58 & \textbf{2.90} \\
		\bottomrule
	\end{tabular}
\end{table}

Table~\ref{tab:final_results} summarizes the validation results across all three samples. The pooled analysis yields a subgroup treatment effect of 0.0448 with a $t$-statistic of 2.90, corresponding to a two-sided $p$-value of 0.004. This highly significant result confirms that the top 10\% of customers identified by our StaDISC analysis experience treatment effects substantially larger than the population average---approximately 4.5 percentage points higher, and in the opposite direction from the negative population ATE.

\section{Discussion}

\subsection{Summary of Findings}
Our application of the StaDISC methodology to the Turkish bank overdraft experiment has yielded a nuanced picture of heterogeneous treatment effects that neither validates nor dismisses the promise of personalized interventions. On one hand, global CATE estimation comprehensively fails: all 23 estimators we trained exhibit poor calibration, with $\rsquare$-scores that are positive on training data but negative on validation folds. This systematic overfitting indicates that the sophisticated machine learning methods commonly used for treatment effect heterogeneity are capturing noise rather than signal when applied to this dataset.

On the other hand, the StaDISC framework reveals that meaningful heterogeneity does exist, but it must be sought in a more circumscribed manner. By focusing on quantile-based top subgroups rather than full CATE predictions, we identify stable pockets of the feature space where treatment effects reliably exceed the population average. The key insight is that even when an estimator cannot accurately predict the magnitude of individual treatment effects, it may still correctly identify which individuals belong to a high-response subgroup.

Among the 23 estimators considered, the R-learner with Random Forest outcome model and Lasso effect learner emerges as the most stable performer. This hybrid architecture appears to balance flexibility and regularization in a way that resists overfitting to particular data partitions. Crucially, this estimator achieves consistently high $t$-statistics across all three CV perturbations---a stringent criterion that most competitors fail to meet.

The final validation confirms the practical significance of our findings. The top 10\% subgroup identified by the selected estimator exhibits a pooled treatment effect of 0.045 with a $t$-statistic of 2.90 ($p < 0.01$). This effect is not only statistically significant but also economically meaningful: it is approximately three times larger than the population average and in the opposite direction. While the average customer shows a small negative response to the bill payment messaging intervention, the top 10\% shows a substantial positive response.

\subsection{Limitations}
This analysis implements the StaDISC methodology through Stage 2 (stability-driven ranking and aggregation) but does not proceed to Stage 3 (CellSearch for interpretable subgroups). As a consequence, the identified subgroup is defined by model predictions rather than by interpretable combinations of customer features. We cannot provide specific targeting rules of the form ``customers with balance greater than X and tenure less than Y.'' The mechanism driving the heterogeneity---why certain customers respond positively while others respond negatively---remains uncharacterized.

This limitation has practical implications for deployment. A bank wishing to target high-responders based on our analysis would need to implement the fitted \texttt{r\_rflasso} model in its production systems and compute predicted treatment effects for each customer. This is more complex than a simple rule-based targeting strategy, though potentially more accurate. Future work should implement the CellSearch procedure to determine whether the quantile-based subgroup can be approximated by interpretable feature combinations that would enable simpler deployment.

\subsection{Implications for Practice}
Our results carry several implications for practitioners considering personalized messaging interventions in consumer finance.

First, global CATE estimates should be treated with substantial skepticism. The strong performance of machine learning methods on training data provides no guarantee of generalization, and the fundamental problem of missing counterfactuals means that standard validation procedures from supervised learning do not directly apply. Calibration checks, such as the $\rsquare$-scores employed here, should precede any deployment of personalized treatment rules.

Second, when heterogeneity signals are weak---as they often are in field experiments powered only to detect average effects---practitioners may achieve more reliable results by targeting the extreme quantiles rather than attempting fine-grained personalization. Our analysis suggests that the top 10--20\% of predicted responders can be identified with reasonable stability, even when the full distribution of treatment effects cannot be accurately estimated.

Third, stability across perturbations is essential for credible heterogeneity claims. A single train-test split provides insufficient evidence, as many estimators achieve impressive performance on particular data partitions only to collapse when the partition changes. Multiple CV splits with different random seeds, along with perturbations to feature definitions and outcome specifications, should be standard practice for any serious investigation of treatment effect heterogeneity.

\section{Conclusion}

Applying the StaDISC framework to the Alan et al.\ Turkish bank overdraft experiment reveals a pattern that is likely common across field experiments in behavioral economics and consumer finance: sophisticated CATE estimators fail to generalize at the global level, but stable local heterogeneity can be extracted through careful calibration checks and stability-driven ranking. The top 10\% of customers identified by an R-learner ensemble exhibit treatment effects approximately three times the population average, with statistical significance confirmed on both training and holdout data.

This analysis demonstrates the value of the PCS framework's emphasis on stability alongside predictability. By demanding that conclusions survive multiple perturbations, StaDISC filters out the ephemeral patterns that plague modern causal machine learning and surfaces effects that are more likely to replicate in future interventions. The methodology provides a principled middle ground between the extremes of ignoring heterogeneity entirely and over-fitting complex CATE models that fail to generalize.

The bill payment messaging intervention studied here illustrates both the promise and the limitations of personalized behavioral interventions in finance. While the average treatment effect is small and negative, our analysis identifies a meaningful subgroup for whom the intervention produces substantial positive effects. Characterizing this subgroup through interpretable features---the natural next step in the StaDISC workflow---would enable targeted deployment that maximizes intervention effectiveness while avoiding the costs of treating non-responders.

\bibliographystyle{plainnat}
\bibliography{references}
	
\end{document}